% ============================================================
%  part7/ch38-lcdm-comparison.tex
%  Chapter 38: Comparison with LCDM
% ============================================================

\chapter{Comparison with $\Lambda$CDM}
\label{ch:lcdm-comparison}

\epigraph{%
  The strongest version of an alternative theory is not\\
  ``why the standard model is wrong''\\
  but ``how two different dynamical pictures\\
  generate similar observables.''%
}{---}

\section{Structure of the comparison}

For each observable, we ask:
\begin{enumerate}[leftmargin=2em]
  \item What is the $\Lambda$CDM mechanism?
  \item What is the RSVP mechanism?
  \item Do they make different predictions?
  \item If so, which current or future observations could
    discriminate?
\end{enumerate}

\section{Observable comparison table}

\begin{center}
\small
\begin{tabular}{@{}p{2.5cm}p{3.5cm}p{3.5cm}p{2.5cm}@{}}
\toprule
\textbf{Observable}
& \textbf{$\Lambda$CDM mechanism}
& \textbf{RSVP mechanism}
& \textbf{Status} \\
\midrule
Hubble expansion
& Metric scale factor $a(t)$
& Path-averaged relaxation $H_{\mathrm{eff}} = c\langle\mathcal{R}\rangle$
& \statusconjectured{} \\
\addlinespace
Redshift
& $1+z = a(t_{\mathrm{obs}})/a(t_{\mathrm{emit}})$
& $1+z = \exp(\int_\gamma\mathcal{R}_{\mathrm{eff}}\,ds)$
& \statusconjectured{} \\
\addlinespace
Large-scale structure
& Gravitational amplification of primordial fluctuations
& Coarsening defect network of quenched $(\Phi,\mathbf{v},S)$
& \statusconjectured{} \\
\addlinespace
Void statistics
& Under-dense regions expand with background
& Lamphrodyne-dominated basins in entropy descent
& \statusconjectured{} \\
\addlinespace
Filament topology
& Gravitational collapse into sheets and filaments
& Defect network of coarsening scalar-vector field
& \statusconjectured{} \\
\addlinespace
BAO peak
& Sound horizon at recombination imprinted in density field
& Characteristic coarsening length $\ell(t_{\mathrm{rec}})$
& \statusopen{} \\
\addlinespace
CMB power spectrum
& Quantum fluctuations during inflation
& Initial conditions for RSVP coarsening (to be specified)
& \statusopen{} \\
\addlinespace
Weak lensing
& Integrated mass along line of sight
& Integrated $\Phi$ and $\ell^2$ along line of sight
& \statusopen{} \\
\addlinespace
Type Ia SN distances
& Luminosity distance via $a(t)$
& Luminosity distance via $\exp(\int\mathcal{R}\,ds)$
& \statusconjectured{} \\
\bottomrule
\end{tabular}
\end{center}

\section{The key distinguishing prediction}

The most distinctive RSVP prediction concerns void lines of sight.

In $\Lambda$CDM: the ISW effect predicts a temperature decrement
(blueshift) for photons traversing voids, because the potential
well is shallower when the photon exits than when it entered.

In RSVP: photons traversing void interiors should accumulate
\emph{less} relaxation than average, producing a systematic
redshift deficit relative to path length, bounded below by the
path-averaged positivity constraint.

\begin{ObsConseq}[Distinguishing prediction]
  \textbf{Prediction:} Void lines of sight show redshift deficit
  relative to equal path length through average density regions.

  \textbf{Sign:} RSVP predicts smaller $z$ through voids;
  $\Lambda$CDM ISW predicts a temperature anisotropy of
  opposite sign.

  \textbf{Current data:} ISW signal from void stacking
  (Planck collaboration, DES, KiDS).

  \textbf{Status:} \statusopen{}
  The quantitative prediction requires a completed eikonal
  derivation of $\mathcal{R}_{\mathrm{eff}}$.
\end{ObsConseq}
